\documentclass{beamer}
\usepackage[utf8]{inputenc}
\usepackage[spanish]{babel}
\usepackage{listings}
\usepackage{xcolor}
\usepackage{tikz}
\usepackage{tcolorbox}
\usepackage{fontawesome5}

% Configuración de colores - Gruvbox Light con marrón principal
\definecolor{bg0}{HTML}{fbf1c7}
\definecolor{bg1}{HTML}{ebdbb2}
\definecolor{bg2}{HTML}{d5c4a1}
\definecolor{bg3}{HTML}{bdae93}
\definecolor{fg0}{HTML}{282828}
\definecolor{fg1}{HTML}{3c3836}
\definecolor{fg2}{HTML}{504945}
\definecolor{aqua}{HTML}{689d6a}
\definecolor{aquabright}{HTML}{427b58}
\definecolor{blue}{HTML}{458588}
\definecolor{red}{HTML}{cc241d}
\definecolor{orange}{HTML}{d65d0e}
\definecolor{yellow}{HTML}{d79921}
\definecolor{green}{HTML}{98971a}
\definecolor{purple}{HTML}{b16286}
\definecolor{brown}{HTML}{AB886D}
\definecolor{brownbright}{HTML}{8B6849}

% Colores para código
\definecolor{codegreen}{HTML}{98971a}
\definecolor{codegray}{HTML}{7c6f64}
\definecolor{codepurple}{HTML}{b16286}
\definecolor{backcolour}{HTML}{f9f5d7}

% Configuración de código
\lstdefinestyle{pythonstyle}{
    backgroundcolor=\color{backcolour},   
    commentstyle=\color{codegreen},
    keywordstyle=\color{brownbright}\bfseries,
    numberstyle=\tiny\color{codegray},
    stringstyle=\color{codepurple},
    basicstyle=\ttfamily\footnotesize\color{fg0},
    breakatwhitespace=false,         
    breaklines=true,                 
    captionpos=b,                    
    keepspaces=true,                 
    numbers=left,                    
    numbersep=3pt,                  
    showspaces=false,                
    showstringspaces=false,
    showtabs=false,                  
    tabsize=2,
    language=Python,
    frame=single,
    rulecolor=\color{brown},
    xleftmargin=8pt,
    xrightmargin=8pt,
    inputencoding=utf8,
    extendedchars=true,
    literate=
        {á}{{\'a}}1 {é}{{\'e}}1 {í}{{\'i}}1 {ó}{{\'o}}1 {ú}{{\'u}}1
        {Á}{{\'A}}1 {É}{{\'E}}1 {Í}{{\'I}}1 {Ó}{{\'O}}1 {Ú}{{\'U}}1
        {ñ}{{\~n}}1 {Ñ}{{\~N}}1
        {¿}{{?`}}1 {¡}{{!`}}1
}

\lstset{style=pythonstyle}

% Tema personalizado con marrón principal
\usetheme{Madrid}
\setbeamercolor{structure}{fg=brown}
\setbeamercolor{palette primary}{bg=brown,fg=bg0}
\setbeamercolor{palette secondary}{bg=brownbright,fg=bg0}
\setbeamercolor{palette tertiary}{bg=fg1,fg=bg0}
\setbeamercolor{palette quaternary}{bg=brown,fg=bg0}
\setbeamercolor{block title}{bg=brown,fg=bg0}
\setbeamercolor{block body}{bg=bg1,fg=fg0}
\setbeamercolor{frametitle}{bg=brown,fg=bg0}
\setbeamercolor{title}{fg=fg0}
\setbeamercolor{subtitle}{fg=fg1}
\setbeamercolor{normal text}{fg=fg0,bg=bg0}
\setbeamercolor{item}{fg=brown}
\setbeamercolor{section in toc}{fg=brown}

% Cajas personalizadas con marrón
\newtcolorbox{conceptbox}[1]{
    colback=bg1,
    colframe=brown,
    fonttitle=\bfseries\small,
    title=#1,
    coltitle=fg0,
    left=3pt,
    right=3pt,
    top=3pt,
    bottom=3pt,
    boxsep=2pt
}

\newtcolorbox{notebox}{
    colback=yellow!15,
    colframe=orange,
    leftrule=2mm,
    rightrule=0mm,
    toprule=0mm,
    bottomrule=0mm,
    arc=0mm,
    left=3pt,
    right=3pt,
    top=2pt,
    bottom=2pt
}

\title{\textbf{Introducción a la Programación con Python}}
\subtitle{Semana 4: Librerías y Módulos}
\author{$\nabla$Nablabs}
\date{}
\begin{document}

\begin{frame}
\titlepage
\end{frame}

\begin{frame}
\frametitle{Contenido}
\tableofcontents
\end{frame}

\section{¿Qué son los módulos y paquetes?}

\begin{frame}{Módulos y Paquetes \faBoxOpen}
\begin{conceptbox}{¿Qué son?}
\small
\begin{itemize}
    \item[\faCube] \textbf{Módulo}: Archivo .py con funciones y variables reutilizables
    \item[\faBoxes] \textbf{Paquete}: Colección de módulos organizados en carpetas
    \item[\faRecycle] Permiten reutilizar código sin repetirlo
    \item[\faUsers] Acceso a código creado por otros programadores
\end{itemize}
\end{conceptbox}

\vspace{0.2cm}
\begin{center}
\large\textcolor{brown}{\faBook\ ¡No reinventes la rueda!}
\end{center}
\end{frame}

\section{Importación de módulos}

\begin{frame}[fragile]{Importar con import \faFileImport}
\begin{conceptbox}{¿Cómo funciona?}
Trae todo el módulo y accedes con punto (.)
\end{conceptbox}

\begin{lstlisting}
import math

# Usar funciones del módulo
raiz = math.sqrt(16)
print(f"Raíz cuadrada: {raiz}")

pi = math.pi
print(f"Valor de pi: {pi}")
\end{lstlisting}

\vspace{-0.2cm}
\begin{tcolorbox}[colback=green!10,colframe=green,title={\small\faArrowCircleRight \ Salida},coltitle=fg0,left=3pt,right=3pt,top=2pt,bottom=2pt,boxsep=1pt]
\small\texttt{Raíz cuadrada: 4.0\\
Valor de pi: 3.141592653589793}
\end{tcolorbox}

\vspace{-0.1cm}
\begin{notebox}
\small\faInfoCircle\ Siempre usa el prefijo del módulo: \texttt{math.sqrt()}
\end{notebox}
\end{frame}

\begin{frame}[fragile]{Importar con alias \faTag}
\begin{conceptbox}{¿Para qué?}
Acorta el nombre del módulo para escribir menos
\end{conceptbox}

\begin{lstlisting}
import math as m
import statistics as stats

raiz = m.sqrt(25)
media = stats.mean([1, 2, 3, 4, 5])

print(f"Raíz: {raiz}")
print(f"Media: {media}")
\end{lstlisting}

\vspace{-0.2cm}
\begin{tcolorbox}[colback=green!10,colframe=green,title={\small\faArrowCircleRight \ Salida},coltitle=fg0,left=3pt,right=3pt,top=2pt,bottom=2pt,boxsep=1pt]
\small\texttt{Raíz: 5.0\\
Media: 3}
\end{tcolorbox}

\vspace{-0.1cm}
\begin{notebox}
\small\faLightbulb\ Útil con nombres largos o convenciones (numpy as np)
\end{notebox}
\end{frame}

\section{Importación selectiva}

\begin{frame}[fragile]{from ... import ... \faFilter}
\begin{conceptbox}{¿Cómo funciona?}
Importa solo funciones específicas sin necesidad del prefijo
\end{conceptbox}

\begin{lstlisting}
from math import sqrt, pi, pow

# Usar directamente sin "math."
raiz = sqrt(16)
potencia = pow(2, 3)

print(f"Raíz: {raiz}")
print(f"Pi: {pi}")
print(f"Potencia: {potencia}")
\end{lstlisting}

\vspace{-0.2cm}
\begin{tcolorbox}[colback=green!10,colframe=green,title={\small\faArrowCircleRight \ Salida},coltitle=fg0,left=3pt,right=3pt,top=2pt,bottom=2pt,boxsep=1pt]
\small\texttt{Raíz: 4.0\\
Pi: 3.141592653589793\\
Potencia: 8.0}
\end{tcolorbox}
\end{frame}

\begin{frame}[fragile]{Importar todo con * \faAsterisk}
\begin{columns}[T]
\column{0.58\textwidth}
\begin{lstlisting}
from math import *

# Todo disponible sin prefijo
resultado = sqrt(9) + pi
print(f"Resultado: {resultado}")
\end{lstlisting}

\vspace{-0.2cm}
\begin{tcolorbox}[colback=green!10,colframe=green,title={\small\faArrowCircleRight \ Salida},coltitle=fg0,left=3pt,right=3pt,top=2pt,bottom=2pt,boxsep=1pt]
\small\texttt{Resultado: 6.141592653589793}
\end{tcolorbox}

\column{0.38\textwidth}
\begin{tcolorbox}[colback=red!10,colframe=red,title={\small\faExclamationTriangle \ Advertencia},coltitle=fg0,left=2pt,right=2pt,top=2pt,bottom=2pt]
\tiny\textbf{No recomendado:}
\begin{itemize}
    \item Contamina el namespace
    \item Dificulta saber origen
    \item Puede causar conflictos
\end{itemize}
\end{tcolorbox}
\end{columns}
\end{frame}

\section{Librería estándar de Python}

\begin{frame}{Librería estándar \faBook}
\begin{conceptbox}{¿Qué es?}
Colección de módulos que vienen incluidos con Python
\end{conceptbox}

\vspace{0.2cm}
\begin{columns}[T]
\column{0.48\textwidth}
\textbf{Módulos comunes:}
\small
\begin{itemize}
    \item[\faDice] \texttt{random} - Aleatoriedad
    \item[\faChartBar] \texttt{statistics} - Estadística
    \item[\faTerminal] \texttt{sys} - Sistema
    \item[\faFolder] \texttt{os} - Sistema operativo
\end{itemize}

\column{0.48\textwidth}
\textbf{Otros útiles:}
\small
\begin{itemize}
    \item[\faClock] \texttt{datetime} - Fechas
    \item[\faFile] \texttt{json} - Formato JSON
    \item[\faExchangeAlt] \texttt{re} - Expresiones regulares
    \item[\faGlobe] \texttt{urllib} - URLs
\end{itemize}
\end{columns}

\vspace{0.3cm}
\begin{center}
\textcolor{brown}{\faInfoCircle\ ¡No necesitas instalar nada!}
\end{center}
\end{frame}

\section{Módulo random}

\begin{frame}[fragile]{random: Números aleatorios \faDice}
\begin{lstlisting}
import random

# Número aleatorio entre 0 y 1
aleatorio = random.random()

# Entero aleatorio en rango
dado = random.randint(1, 6)

# Elegir elemento de lista
frutas = ["manzana", "banana", "cereza"]
fruta = random.choice(frutas)

print(f"Aleatorio: {aleatorio:.2f}")
print(f"Dado: {dado}")
print(f"Fruta: {fruta}")
\end{lstlisting}

\vspace{-0.2cm}
\begin{tcolorbox}[colback=green!10,colframe=green,title={\small\faArrowCircleRight \ Salida (ejemplo)},coltitle=fg0,left=3pt,right=3pt,top=2pt,bottom=2pt,boxsep=1pt]
\small\texttt{Aleatorio: 0.73\\
Dado: 4\\
Fruta: banana}
\end{tcolorbox}
\end{frame}

\begin{frame}[fragile]{random: Más funciones \faShuffle}
\begin{lstlisting}
import random

# Mezclar lista
numeros = [1, 2, 3, 4, 5]
random.shuffle(numeros)
print(f"Mezclados: {numeros}")

# Seleccionar múltiples elementos
colores = ["rojo", "azul", "verde", "amarillo"]
seleccion = random.sample(colores, 2)
print(f"Selección: {seleccion}")
\end{lstlisting}

\vspace{-0.2cm}
\begin{tcolorbox}[colback=green!10,colframe=green,title={\small\faArrowCircleRight \ Salida (ejemplo)},coltitle=fg0,left=3pt,right=3pt,top=2pt,bottom=2pt,boxsep=1pt]
\small\texttt{Mezclados: [3, 1, 5, 2, 4]\\
Selección: ['verde', 'rojo']}
\end{tcolorbox}

\vspace{-0.1cm}
\begin{notebox}
\small\faExclamationTriangle\ \texttt{shuffle()} modifica la lista original
\end{notebox}
\end{frame}

\section{Módulo statistics}

\begin{frame}[fragile]{statistics: Estadística básica \faChartBar}
\begin{lstlisting}
import statistics as stats

datos = [10, 20, 30, 40, 50]

media = stats.mean(datos)
mediana = stats.median(datos)
desv_std = stats.stdev(datos)

print(f"Media: {media}")
print(f"Mediana: {mediana}")
print(f"Desviación estándar: {desv_std:.2f}")
\end{lstlisting}

\vspace{-0.2cm}
\begin{tcolorbox}[colback=green!10,colframe=green,title={\small\faArrowCircleRight \ Salida},coltitle=fg0,left=3pt,right=3pt,top=2pt,bottom=2pt,boxsep=1pt]
\small\texttt{Media: 30.0\\
Mediana: 30\\
Desviación estándar: 15.81}
\end{tcolorbox}

\vspace{-0.1cm}
\begin{notebox}
\small\faLightbulb\ Ideal para análisis rápidos de datos
\end{notebox}
\end{frame}

\section{Módulos sys y os}

\begin{frame}[fragile]{sys: Información del sistema \faTerminal}
\begin{lstlisting}
import sys

# Versión de Python
print(f"Versión: {sys.version}")

# Argumentos de línea de comandos
print(f"Argumentos: {sys.argv}")

# Plataforma
print(f"Plataforma: {sys.platform}")

# Salir del programa
# sys.exit()
\end{lstlisting}

\vspace{-0.2cm}
\begin{tcolorbox}[colback=green!10,colframe=green,title={\small\faArrowCircleRight \ Salida (ejemplo)},coltitle=fg0,left=3pt,right=3pt,top=2pt,bottom=2pt,boxsep=1pt]
\small\texttt{Versión: 3.11.0\\
Argumentos: ['programa.py']\\
Plataforma: win32}
\end{tcolorbox}
\end{frame}

\begin{frame}[fragile]{os: Sistema operativo \faFolder}
\begin{lstlisting}
import os

# Directorio actual
directorio = os.getcwd()
print(f"Directorio: {directorio}")

# Listar archivos
archivos = os.listdir(".")
print(f"Archivos: {archivos[:3]}")  # Primeros 3

# Crear directorio
# os.mkdir("nueva_carpeta")

# Verificar si existe
existe = os.path.exists("archivo.txt")
print(f"Existe: {existe}")
\end{lstlisting}

\vspace{-0.2cm}
\begin{tcolorbox}[colback=green!10,colframe=green,title={\small\faArrowCircleRight \ Salida (ejemplo)},coltitle=fg0,left=3pt,right=3pt,top=2pt,bottom=2pt,boxsep=1pt]
\small\texttt{Directorio: /home/usuario\\
Archivos: ['programa.py', 'datos.txt', ...]\\
Existe: False}
\end{tcolorbox}
\end{frame}

\section{Uso de pip}

\begin{frame}[fragile]{pip: Gestor de paquetes \faDownload}
\begin{conceptbox}{¿Qué es pip?}
Herramienta para instalar paquetes externos de PyPI (Python Package Index)
\end{conceptbox}

\vspace{0.2cm}
\textbf{Comandos básicos en terminal:}
\begin{lstlisting}[language=bash, numbers=none]
# Instalar paquete
pip install requests
# Instalar versión específica
pip install numpy==1.24.0
# Desinstalar
pip uninstall requests
# Listar instalados
pip list
# Actualizar paquete
pip install --upgrade requests
\end{lstlisting}
\end{frame}

\begin{frame}[fragile]{Ejemplo: Usar paquete instalado \faRocket}
\begin{lstlisting}
# Primero instalar: pip install requests
import requests

# Hacer petición HTTP
respuesta = requests.get("https://api.github.com")

print(f"Código: {respuesta.status_code}")
print(f"Tipo: {respuesta.headers['Content-Type']}")
\end{lstlisting}

\vspace{-0.2cm}
\begin{tcolorbox}[colback=green!10,colframe=green,title={\small\faArrowCircleRight \ Salida},coltitle=fg0,left=3pt,right=3pt,top=2pt,bottom=2pt,boxsep=1pt]
\small\texttt{Código: 200\\
Tipo: application/json; charset=utf-8}
\end{tcolorbox}

\vspace{-0.1cm}
\begin{notebox}
\small\faInfoCircle\ Más de 400,000 paquetes disponibles en PyPI
\end{notebox}
\end{frame}

\section{Creación de módulos propios}

\begin{frame}[fragile]{Crear tu propio módulo \faWrench}
\begin{conceptbox}{¿Cómo crear un módulo?}
Simplemente crea un archivo .py con funciones
\end{conceptbox}

\textbf{Archivo:} \texttt{matematicas.py}
\begin{lstlisting}
def sumar(a, b):
    """Suma dos números"""
    return a + b

def multiplicar(a, b):
    """Multiplica dos números"""
    return a * b

PI = 3.14159
\end{lstlisting}

\vspace{-0.1cm}
\begin{notebox}
\small\faLightbulb\ El nombre del archivo es el nombre del módulo
\end{notebox}
\end{frame}

\begin{frame}[fragile]{Usar tu módulo \faPlay}
\textbf{Archivo:} \texttt{programa.py} (en la misma carpeta)
\begin{lstlisting}
import matematicas

resultado = matematicas.sumar(5, 3)
print(f"Suma: {resultado}")

producto = matematicas.multiplicar(4, 6)
print(f"Producto: {producto}")

print(f"Pi: {matematicas.PI}")
\end{lstlisting}

\vspace{-0.2cm}
\begin{tcolorbox}[colback=green!10,colframe=green,title={\small\faArrowCircleRight \ Salida},coltitle=fg0,left=3pt,right=3pt,top=2pt,bottom=2pt,boxsep=1pt]
\small\texttt{Suma: 8\\
Producto: 24\\
Pi: 3.14159}
\end{tcolorbox}

\vspace{-0.1cm}
\begin{notebox}
\small\faCheckCircle\ ¡Así de fácil es crear y usar tus módulos!
\end{notebox}
\end{frame}

\begin{frame}[fragile]{Módulo con if \_\_name\_\_ == "\_\_main\_\_" \faCodeBranch}
\textbf{Archivo:} \texttt{calculadora.py}
\begin{lstlisting}
def sumar(a, b):
    return a + b

def restar(a, b):
    return a - b

# Solo se ejecuta si se corre directamente
if __name__ == "__main__":
    print("Probando calculadora...")
    print(f"5 + 3 = {sumar(5, 3)}")
    print(f"5 - 3 = {restar(5, 3)}")
\end{lstlisting}

\vspace{-0.1cm}
\begin{notebox}
\small\faInfoCircle\ El bloque no se ejecuta cuando importas el módulo
\end{notebox}
\end{frame}

\section{Documentación y búsqueda}

\begin{frame}{Buscar y documentarse \faSearch}
\begin{conceptbox}{¿Dónde buscar librerías?}
\small
\begin{itemize}
    \item[\faGlobe] \textbf{PyPI}: \texttt{pypi.org} - Repositorio oficial
    \item[\faBook] \textbf{Docs Python}: \texttt{docs.python.org} - Librería estándar
    \item[\faGithub] \textbf{GitHub}: Código fuente y ejemplos
    \item[\faStackOverflow] \textbf{Stack Overflow}: Preguntas y respuestas
\end{itemize}
\end{conceptbox}

\vspace{0.2cm}

\end{frame}

\begin{frame}{Librerías populares \faStar}
\begin{conceptbox}{Más utilizadas en Python}
\small
\begin{itemize}
    \item[\faChartLine] \textbf{NumPy}: Computación numérica
    \item[\faTable] \textbf{Pandas}: Análisis de datos
    \item[\faChartPie] \textbf{Matplotlib}: Visualización de datos
    \item[\faGlobe] \textbf{Requests}: Peticiones HTTP
    \item[\faFlask] \textbf{Flask/Django}: Desarrollo web
    \item[\faRobot] \textbf{Scikit-learn}: Machine Learning
    \item[\faBrain] \textbf{TensorFlow/PyTorch}: Deep Learning
    \item[\faFileAlt] \textbf{Beautiful Soup}: Web scraping
\end{itemize}
\end{conceptbox}
\end{frame}

\section{Ejemplo integrador}

\begin{frame}[fragile]{Proyecto: Analizador de notas \faGraduationCap}
\textbf{Archivo:} \texttt{notas\_utils.py}
\begin{lstlisting}
import statistics

def calcular_promedio(notas):
    """Calcula el promedio de las notas"""
    return statistics.mean(notas)

def nota_maxima(notas):
    """Retorna la nota máxima"""
    return max(notas)

def nota_minima(notas):
    """Retorna la nota mínima"""
    return min(notas)

def estado(promedio):
    """Determina si aprobó o reprobó"""
    return "Aprobado" if promedio >= 60 else "Reprobado"
\end{lstlisting}
\end{frame}

\begin{frame}[fragile]{Proyecto: Programa principal \faDesktop}
\textbf{Archivo:} \texttt{analizar.py}
\begin{lstlisting}
import notas_utils as nu
import random
# Generar notas aleatorias para demostración
notas = [random.randint(40, 100) for _ in range(5)]
print(f"Notas: {notas}")
print(f"Promedio: {nu.calcular_promedio(notas):.1f}")
print(f"Máxima: {nu.nota_maxima(notas)}")
print(f"Mínima: {nu.nota_minima(notas)}")
promedio = nu.calcular_promedio(notas)
print(f"Estado: {nu.estado(promedio)}")
\end{lstlisting}

\vspace{-0.2cm}
\begin{tcolorbox}[colback=green!10,colframe=green,title={\small\faArrowCircleRight \ Salida (ejemplo)},coltitle=fg0,left=3pt,right=3pt,top=2pt,bottom=2pt,boxsep=1pt]
\small\texttt{Notas: [85, 72, 91, 68, 77]\\
Promedio: 78.6\\
Máxima: 91\\
Mínima: 68\\
Estado: Aprobado}
\end{tcolorbox}
\end{frame}

\begin{frame}{Resumen \faListUl}
\begin{conceptbox}{¿Qué aprendimos hoy?}
\small
\begin{itemize}
    \item[\faCheck] Módulos y paquetes
    \item[\faCheck] Importar con \texttt{import} y \texttt{from ... import}
    \item[\faCheck] Librería estándar de Python
    \item[\faCheck] Módulos: random, statistics, sys, os
    \item[\faCheck] Uso de pip para instalar paquetes
    \item[\faCheck] Creación de módulos propios
    \item[\faCheck] Documentación y búsqueda de librerías
\end{itemize}
\end{conceptbox}

\vspace{0.3cm}
\begin{center}
\large\textcolor{brown}{\faBoxOpen\ ¡Explora el ecosistema Python!}
\end{center}
\end{frame}

\begin{frame}
\begin{center}
\Huge\textcolor{brown}{\faQuestionCircle}\\
\vspace{0.5cm}
\Huge ¡Preguntas!\\
\vspace{1cm}
\Large\textcolor{brownbright}{Gracias por su atención}
\end{center}
\end{frame}

\end{document}
